\documentclass[12pt,a4paper]{article}

\usepackage[T1]{fontenc}
\usepackage[utf8]{inputenc}
\usepackage[english]{babel}
\usepackage{lmodern}
\usepackage{microtype}
\usepackage{geometry}

\geometry{margin=2.5cm}
\setlength{\parindent}{0pt}
\setlength{\parskip}{0.7em}
\pagestyle{empty}

\begin{document}

\section*{Summary}

This thesis examines fine-tuning small open-weight language models for binary
email classification as legitimate or unwanted. After reviewing the available
models, Qwen3-0.6B was selected for the experiments. Low-rank adaptation
(LoRA) was used. Sixteen configurations and model states
saved at successive training stages were compared. Four ways of using the
model as a classifier were then evaluated. The comparison also included the
unmodified model and six other classifiers. The final evaluation used four
disjoint test sets containing 2,000 messages each. Across the four test sets,
the best fine-tuned variant achieved a mean classification accuracy of 97.8\%.
DistilBERT, the most accurate of the other classifiers, achieved a mean
accuracy of 94.4\% and processed messages approximately 2.3 times faster than
Qwen3.
Using the fine-tuned model in an email filter requires accounting for the
computational cost of classification and evaluating it on current email
traffic.

\end{document}
